\documentclass[openany,fontset=fandol]{ctexbook}
\usepackage{amsmath,amssymb}
\usepackage[unicode,hidelinks]{hyperref}
\newcommand{\note}[1]{\unskip\nobreak\footnote{译注：#1}}
\begin{document}
\title{代数几何基础第五卷预备笔记\\第15节：推广计划结语}
\author{据用户上传英译编辑稿译出}
\date{2026年11月19日}
\maketitle
\chapter*{范围与校订状态}
本批译出底本第61页最后两个完整段落，从“The passage in No. 4 from a generic section to a general section”至“if you have special difficulties”。至此，本轮推进至§15末。

下一接续位置为底本第62页“Part II [Tr]”、§16“New EGA V”及“New \# Singular and supersingular set of a function and differential criteria”。本批没有译入§16。

保留7.3后的撇号、“quant (???)”及关于8.7(c)的句法疑点。第3、6节推广仍按作者原话处于等待状态，未宣称全部公理化或旧缺文已经完成。独立续译，未并主稿。
\chapter*{§15 前述若干结果的公理化（续）}
第4节中从泛截面到一般截面的过渡，以及第5节的论述，在当前背景下也有效（不过，到这里，它们已成为同义反复，或只是对第8、9、12节的重新表述；我们对于以这种形式陈述它们仍有所犹豫）。\note{底本有“evidently valid”，此处按原句结论译为“也有效”，不将其当作本译者已补全的一般化证明。“paragraphs 8, 9, 12”沿用本稿编号体系译为第8、9、12节。}
同样，若$k$代数闭，7.1的论述在所有情形中都有效；若再假设$G$在$k$上有理，那么即使只要求$k$无限也成立。特殊情形7.2、7.3$'$也是如此；至于〔quant，???〕7.4的结果，则是针对超平面截面情形的特殊应用。如我已经说过的，第3节和第6节仍须暂缓，等待Zariski定理的推广。\note{底本原句为“and the special cases 7.2, 7.3' quant (???) the result 7.4”。正文保留7.3后的撇号和“quant”旁的问号，仅按句意以“至于”作暂定连接，不把编号改成另一个确定条目。“rational over $k$”译为“在$k$上有理”，不是指$G$仅有一个$k$有理点；括号“suspended pending”明确保留为等待推广的状态。}

还须推广第8节的结果（第12节曾重新考虑过），它们由此呈现出更令人满意的面貌。我建议你先在这个背景下表述这些结果，并尝试沿这条路尽可能推进。我觉得，我们至少应当能够重新得到那些不是8.7(c)直接推论的结果（甚至可以尝试提炼出一些公理化条件，使你能够完成8.7的一个变式）。\note{底本此句实际为“recover at least that is not a direct consequence of 8.7 c)”，主语或先行词不完整，正文补出最低限度的“那些……结果”以成句；其中“not”照译为“不是”，没有擅自删去否定词。底本括号在下一句前未闭合，正文按语义在“变式”后补右括号。原文究竟意欲恢复哪些结果及否定词是否可靠仍待核，不据此修复8.7(c)的历史缺文。}
我这里只提出这些建议；不过，若你遇到特殊困难，我愿意再回到这些问题上，给出更多细节。\note{这里的“我”和“你”是原稿作者与收信人之间的叙述，保留建议、印象及愿意进一步说明的语气，不把它改写成本译者已经执行的研究或承诺。}
\end{document}
